%% cover_letter.tex — Cover letter for Cognitive Systems Research submission
\documentclass[12pt,a4paper]{letter}

\usepackage[british]{babel}
\usepackage{microtype}
\usepackage{geometry}
\geometry{margin=2.5cm}

\begin{document}

\begin{letter}{%
  The Editors\\
  \textit{Cognitive Systems Research}\\
  Elsevier B.V.}

\opening{Dear Editors,}

This letter accompanies the submission of the manuscript titled
``\textit{Capability-Asymmetric Multi-Agent Debate Degrades Strong LLM
Reasoners: A Bias-Theoretic Account}'' for consideration in
\textit{Cognitive Systems Research}.

\textbf{What the paper does.}
The paper reports a controlled five-condition experiment in which a strong
large language model (LLM) focal agent debates under increasing numbers of
deliberately weak peers, measuring how peer composition affects final-round
accuracy, correct-to-incorrect flip rates, and the Asch conformity index.
The primary focal agent is \textit{DeepSeek-v4-flash}; a cross-validation
uses \textit{GPT-4o-mini}. Weak peers are \textit{Llama~3.1~8B Instruct}
and \textit{Gemma~3~4B Instruct}, achieving around 30\% accuracy on the
same 300-question pool drawn from MMLU-Pro and GSM8K.

\textbf{Why it belongs in this journal.}
\textit{Cognitive Systems Research} covers computational models of human and
machine cognition, including bias, social influence, and collective
reasoning.
This work brings a classic social psychology paradigm (Asch conformity) to
bear on multi-agent LLM systems, providing the first controlled dose-response
measurement of how capability asymmetry drives conformity-induced errors.
The bias-theoretic framing, the mechanistic decomposition of sycophancy and
bandwagon effects, and the practical mitigation analysis are all directly
relevant to the journal's scope.

\textbf{Key findings.}
The correct-to-incorrect flip rate rises monotonically with dumb peer count
(5.0\% in C2 to 6.5\% in C4) even when the net accuracy outcome is positive
for the stronger model.
For \textit{GPT-4o-mini}, accuracy falls 2.0 percentage points in C3 and
the Asch conformity index reaches 0.172 in C4.
Confidence-weighted peer filtering (C5) fails completely because weak agents
report high confidence in 99.5\% of trials, including incorrect responses.
The McNemar test confirms significant question-level regression under maximal
dumb-peer exposure ($p = 0.0025$, Bonferroni-corrected).

\textbf{Novelty and significance.}
This work differs from prior MAD studies by (a)~independently varying the
count of weak peers in a factorial design, (b)~mapping outcomes explicitly
onto the sycophancy and bandwagon mechanisms through the Asch index, and
(c)~replicating across two focal models of different capability.
To the best of the author's knowledge, no prior paper has subjected a
confidence-weighted peer filter to a deliberate worst-case test of this kind.
The negative result for C5 is a timely and practically important finding.

\textbf{Declarations.}
The manuscript has not been published elsewhere and is not under
consideration by any other journal.
The experiment involves only publicly available language models and benchmark
datasets; no human subjects were involved.
Data availability is described in a dedicated statement included with the
submission.

This paper has been prepared for single-blind review and the author looks
forward to the reviewers' feedback.

\closing{Yours sincerely,}

\vspace{1em}
Koushik Deb\\
{[Department, Institution]}\\
{[Email]}\\
{[Date]}

\end{letter}

\end{document}
